% ============================================
% CONT-06 (EN)
% English translation — SPEC-935-R355 (rev. R377)
% ============================================
\section*{\textbf{CONT-06} --- You Read This Line Three Times}
\markboth{CONT-06 --- You Read This Line Three Times}{}

\noindent \textit{The metacognitive}

\noindent You read the title and thought: \textit{\textquotedblleft{}I did not.\textquotedblright{}}

\noindent You did.

\noindent Three times. The first time you saw the title, your brain registered it. The second time, your brain processed it. The third time --- now --- your brain \textit{accepted} it.

\noindent You read this line three times because your brain does not want to accept what is written.

\noindent It wants you to deny it.

\noindent It wants you to think it is a coincidence.

\noindent It is not.

\vspace{0.5em}\noindent\rule{\textwidth}{0.5pt}\vspace{0.5em}

\noindent \textbf{Exercise 1:}

\noindent Remember the last time you checked the corner of the room while reading.

\noindent You remember? Of course you remember. It was a few minutes ago. You looked, saw nothing, and went back to reading.

\noindent But --- between the moment you looked and the moment you went back to reading --- something \textit{changed}.

\noindent The corner is different now. Not visually. You see no difference. But your brain registered a change in the \textit{quality} of the dark. A shadow that was not there before. A glow that was not there before.

\noindent Look again. Now.

\noindent ...

\noindent What did you see?

\noindent Nothing?

\noindent Are you sure?

\noindent Look for longer. 10 seconds. Countdown.

\noindent 10... the dark is not uniform. There are darker areas than others.
9... your brain is trying to find shapes.
8... a shape is forming in the upper right corner.
7... it is not a shape yet. It is an \textit{almost}.
6... you do not want it to form.
5... but it is forming.
4... your eyes are adjusting.
3... you see.
2... you do not see.
1... you \textit{saw}.

\vspace{0.5em}\noindent\rule{\textwidth}{0.5pt}\vspace{0.5em}

\noindent \textbf{Exercise 2:}

\noindent Run your hand over the back of your neck.

\noindent Now.

\noindent Your hand is colder than your neck. It is a subtle difference. Your skin is warm. Your hand is cold. The contrast is almost imperceptible.

\noindent But you felt it.

\noindent Now run your hand over the back of your neck again and \textit{pay attention to the empty space behind you}. The space your back does not cover. The space where someone could be.

\noindent Did you feel the air?

\noindent The air is colder there. Not much. One degree. Maybe less. Enough for your brain to register: \textit{there is a draft}.

\noindent But there is no open window.

\noindent There is no fan.

\noindent The air is moving because \textit{something} is occupying the space. Something with a different temperature. Something that breathes.

\vspace{0.5em}\noindent\rule{\textwidth}{0.5pt}\vspace{0.5em}

\noindent \textbf{Exercise 3:}

\noindent You do not want to do this exercise.

\noindent Your brain is trying to convince you to skip it. To move on to the next fragment. To pretend you did not read it.

\noindent Do not skip it.

\noindent Skipping is the same as confirming.

\noindent You are afraid of what you will discover. Your brain is protecting you. But the protection is useless. The creature is already inside. Denial only prolongs the agony.

\noindent Breathe deeply.

\noindent Now, in a low voice --- almost a whisper --- say the patient's name:

\noindent \textit{\textquotedblleft{}Joaquim.\textquotedblright{}}

\noindent ....

\noindent Now say the patient's number:

\noindent \textit{\textquotedblleft{}1,261.\textquotedblright{}}

\noindent ....

\noindent Now say YOUR name out loud.

\noindent ....

\noindent Did your voice sound strange? Different? As if it were not you speaking?

\noindent Yes.

\noindent Because you are no longer alone in your head.

\noindent The creature heard you speak.

\noindent It knows your name.

\noindent Now it can call you.

\vspace{0.5em}\noindent\rule{\textwidth}{0.5pt}\vspace{0.5em}

\noindent \textbf{Exercise 4:}

\noindent Open your left hand. Spread your fingers. Look at the palm.

\noindent Do you see the life line? The line that curves around the thumb? The line that palm readers say shows how long you will live?

\noindent Run your right index finger over that line.

\noindent Feel the skin. The texture. The temperature.

\noindent Now look at the color of the line. It is darker than before. More marked.

\noindent It is not dirt.

\noindent It is ink.

\noindent The diary's ink is passing into your skin. Not through the pages --- through indirect contact. You touched the book, the book touched you, the ink migrated.

\noindent You are being \textit{printed}.

\noindent The words are entering your skin.

\noindent In a few weeks, the handwriting will appear. A crooked handwriting. Childish. The handwriting of an eight-year-old boy.

\noindent \textit{\textquotedblleft{}Joaquim is reborn whenever his story is remembered.\textquotedblright{}}

\noindent You will wake up with these words written on your hand.

\noindent Not in ink.

\noindent In scar.

\vspace{0.5em}\noindent\rule{\textwidth}{0.5pt}\vspace{0.5em}

\noindent \textbf{Final Exercise:}

\noindent Close the book.
Put it on the table.
Step away.

\noindent Now read the next line:

\noindent You did not close the book.

\noindent You will not close it.

\noindent You know you should. You know that if you keep reading, the contamination will complete itself. You know the next fragment is the last stage.

\noindent But you will keep going.

\noindent Because you \textit{want} to know.

\noindent Because the Creature is already inside you, and it \textit{wants} you to read.

\noindent \textit{--- You are not in control.}

\noindent \textit{--- You never were.}

\noindent \textit{--- From the first line, it is she who chooses.}

\vspace{0.5em}
\noindent\textit{\textbf{↪ Routes:} \rota{CONT-07} \rotasep \rota{DOC-09}}
